% \subsection{Autoregressive Random Effects Models}

One model that has been developed to analyze RCS data with autocorrelated errors is the Autoregressive Random Effects (ARE) Model \citep{modugno2015multilevel}. Treating observations as the sum of a fixed effects model and an AR(1) time series, ARE models have two levels: a fixed-effects level and a time series level.
% \begin{equation}
%     \label{eq:modugno_level1}
%     y_{ti} = \mathbf{x}_{ti}\beta + u_t + \epsilon_{ti}, \hspace{.4in}  i = 1,...,n, \hspace{.4 in} \epsilon_{ti} \sim \text{i.i.d. } N(0,\sigma_\epsilon^2) 
% \end{equation}

% \begin{equation}
%     \label{eq:modugno_level2}
%     u_t - \phi u_{t-1} = \eta_t + \theta \eta_{t-1} , \hspace{.4in} t=1,...T, \hspace{.4in} \eta_t \sim N(0,\sigma_\eta^2) 
% \end{equation}


Setting $N:=nT$, the total number of observations in the dataset,  we define the fixed effect and temporal fluctuation design matrices $X_f$ and $X_t$ (resp.) to be:


With these matrices we can concisely represent the model governing the vector of observations $\yy$:

$$ \yy = X_f\beta + X_t\uu + \bm{\varepsilon}$$

The method employed to fit such a model is a recursive process that splits the log likelihood function of the model into two terms and then utilizes an EM algorithm for parameter estimates. We introduce some necessary notation for expressing the likelihood function of the full set of parameters:







This conditional splitting is convenient, as $\uu \sim N(\mathbf{0},\GG)$ and $\yy|\uu \sim N(X_f\beta + \uu, \sigma_\varepsilon^2I_N)$. Furthermore, the choice of imposing an AR(1) structure on $\uu$ lends itself to a concise representation of the determinant and inverse of $\GG$:



In the absence of an explicit observation of $\uu$, its estimate is repeatedly calculated and used to update parameter estimates. Its estimation is therefore a key part of this algorithm, and is given by its expectation, conditioned on $\yy$: 


% \begin{equation}
%     \label{cov_uy}
%     \text{Cov}(\uu,\yy) = \GG X_t'
% \end{equation}

% With these pieces we can give the following expression for the expectation of $\uu$, conditioned on $\yy$:

With this estimate now defined, we summarize the fitting of an ARE model:





% This method preserves individual information, but accounts for the possible presence of temporal correlation by adding the time series term. The method is a recursive process with the following steps:



\begin{itemize}
    \item[1)] Calculate initial parameter estimates; 
    \item[2)] Estimate $\uu$ with $\mathbb{E}(\uu|\yy)$;
    \item[3)] Apply $\hat{\uu}$ with an AR(1) likelihood to update $\hat{\xi}_2$ and fit $\yy - X_t\hat{\uu}$ with a fixed effects model to update $\hat{\xi}_1$;
    \item[4)] Calculate and store $l_1(\hat{\xi}_1)$ and $l_2(\hat{\xi}_2)$;
    \item[5)] Repeat steps 2-4 until convergence criteria are met.
\end{itemize}

% All estimate updates are based on explicitly calculable maximum likelihood estimators, with the exception of the AR parameter $\phi$, which is updated via one iteration of the Newton-Raphson scheme since $l_2$ is not linear with respect to $\phi$. In the next section, we examine some potential issues with fitting models with the ARE approach.

